%%% Таблицы: общие настройки оформления %%%

%%% Базовый размер шрифта в таблицах
\AtBeginEnvironment{tabularx}{%
    \footnotesize
    \rowcolors{2}{gray!10}{white}%
}

\AtBeginEnvironment{longtable}{%
    \footnotesize
    \rowcolors{2}{white}{gray!10}%
}

%%% Отступы longtable до и после таблицы
\setlength{\LTpre}{0pt}
\setlength{\LTpost}{0pt}

%%% Вертикальный сдвиг нижней подписи longtable
\newcommand{\LongTableBottomCaptionRaise}[1]{%
    \raisebox{7pt}[0pt][0pt]{#1}%
}

%%% Нижняя подпись для longtable
\newcommand{\longcaption}[1]{%
    \rowcolor{white}%
    \caption*{\LongTableBottomCaptionRaise{#1}}%
}

%%% Пользовательские типы колонок
\newcolumntype{X}{>{\RaggedRight\arraybackslash}p{\dimexpr\linewidth/3-2\tabcolsep\relax}}
\newcolumntype{Y}{>{\RaggedRight\arraybackslash}X}
\newcolumntype{Z}{>{\centering\arraybackslash}X}
\newcolumntype{W}{>{\RaggedLeft\arraybackslash}X}

%%% Единый стиль ячеек заголовка
\newcommand{\TableHeader}[1]{%
    \textcolor{white}{\textbf{#1}}%
}

%%% Белая строка для служебных элементов longtable
\newcommand{\WhiteRow}{%
    \rowcolor{white}%
}

%%% Повторяемая подпись longtable на следующих страницах
\newcommand{\LongTableContinuationCaption}{%
    \WhiteRow
    \caption[]{Продолжение} \tabularnewline
}

%%% Универсальная шапка longtable
\newcommand{\LongTableHeader}[1]{%
    #1
    \endfirsthead
    \LongTableContinuationCaption
    #1
    \endhead
    \hline
    \endfoot
    \endlastfoot
}

%%% Нижняя подпись longtable на каждой странице
\newcommand{\longcaptioneachpage}[1]{%
    \longcaption{#1}
    \endfoot
    \longcaption{#1}
    \endlastfoot
}